\documentclass[../main.tex]{subfiles}

\begin{document}

Data collected during the experiment were analyzed to test our 8 hypotheses (H1-H8), 
which we categorized into three logical clusters based on the software development lifecycle.

\subsection{Learning Phase and Cognitive Integrity (H1, H2, H6)}
\begin{itemize}
    \item \textbf{H1-H2 (Cognitive Erosion)}: The data strongly supported the hypothesis that "vibe coding" causes learners to skip line-by-line interpretation of their code, 
    fundamentally hindering deep learning. Participants who fully delegated coding tasks to the AI bypassed the essential cognitive struggle required to form new skills, 
    resulting in impaired conceptual understanding, diminished code-reading abilities, and a significant reduction in overall knowledge retention. 
    Those employing an "AI Delegation" interaction pattern demonstrated the lowest levels of independent thinking and cognitive engagement. \cite{shen2026how}
    \item \textbf{H6 (The Role of Prior Knowledge)}: Conversely, developers with stronger prior manual programming experience proved much faster at recognizing the AI's limitations. 
    Higher-competence programmers evaluate AI-generated code more critically—actively reviewing, debugging, and striving to understand it. Rather than relying on the AI to 
    generate complete solutions, these experienced students used the tool as a "lifesaver" for conceptual inquiry and clarification, which preserved their learning outcomes 
    and protected them from cognitive erosion. \cite{kusper2025we} \cite{shen2026how}
\end{itemize}

I. Tanulási fázis és Kognitív integritás (H1, H2, H6)
\begin{itemize}
    \item H1-H2 (Kognitív erózió): Az adatok várhatóan alátámasztják, hogy a „vibe coding” során elmarad a kód soronkénti értelmezése, ami akadályozza a mély tanulást.
    \item H6 (Alaptudás szerepe): Azok a fejlesztők, akik rendelkeznek korábbi manuális tapasztalattal, hamarabb felismerik az AI korlátait (mentőöv szerep).
\end{itemize}

II. Hibakeresési fázis és Interakciós dinamika (H3, H4)
\begin{itemize}
    \item H3-H4 (A Debugolási Csapda): A kísérleti csoport hajlamos a „prompt-loop”-ra (ugyanazt a hibát többször visszadobja az AI-nak), ahelyett, hogy elolvasná a hibaüzenetet (manuális kódolvasás elkerülése).
    \item Metakognitív monitoring: A résztvevők tévesen ítélik meg saját haladásukat a generált kód mennyisége alapján.
\end{itemize}

III. Kompetencia-szakadék és Rendszerszemlélet (H5, H7, H8)
\begin{itemize}
    \item H5 (Architekturális vakság): A kód „patchwork” jellege miatt a fejlesztő nem látja át a modulok közötti függőségeket.
    \item H7-H8 (Over-reliance): A túlzott hagyatkozás korrelál a kritikai gondolkodás és a logikai következtetés (inference) csökkenésével.
\end{itemize}

\end{document}
